\section{从 RNN 走向 Attention：为什么要让信息直接相遇？}
\label{intro:21b-why-attention}

我们现在已经有了 Token 向量，也知道 RNN 如何逐个处理位置。但理解 Transformer 不能只背诵“RNN 不好，Attention 更好”。先看看两种模型究竟怎样让远距离的信息发生联系。

\subsection{想解决什么问题：一个指代关系}

考虑一句话：
\begin{center}
\emph{“小明把书放进书包，因为\textbf{它}很重。”}
\end{center}
其中“它”指什么，依赖上下文；把词逐个独立分类可能很困难。

RNN 可以逐步带着隐藏状态往后读。这样早期信息有机会影响后面的词，但信息必须经过多个更新步骤。我们能不能允许当前位置\textbf{直接查看前面相关的位置}？

\subsection{先设计一种没有 Key/Value 名字的直觉方案}

假设模型需要从三个位置读取信息。我们给三个位置各算一个相关性分数：
\[
s=(2,0,1).
\]
直接按分数加权会遇到问题：分数可能是负数、大小也可能没有统一标准。于是先把它们经过 Softmax：
\[
a_i=\frac{e^{s_i}}{e^2+e^0+e^1}.
\]
数值近似为
\[
\vect a\approx(0.665,0.090,0.245).
\]
它们都是非负数，总和约为 $1$。

如果三个位置可以提供二维信息：
\[
\vect v_1=(1,0),\quad
\vect v_2=(0,1),\quad
\vect v_3=(1,1),
\]
那么当前位置读到的信息是
\begin{align*}
\vect y&=a_1\vect v_1+a_2\vect v_2+a_3\vect v_3\\
&\approx(0.910,0.335).
\end{align*}

\textbf{此处的核心创新不是指数函数本身，而是“根据当前位置对各信息源的相关性，动态分配读取权重”。}这套设计仍缺一个问题：相关性分数从哪里来？

\subsection{Query、Key、Value：给计算中的三种角色命名}

我们需要一个可以训练的匹配机制。给当前查询位置一个\textbf{Query} $\vect q$，给每个候选位置一个\textbf{Key} $\vect k_i$，让相似度分数由点积产生：
\begin{equation}
s_i=\vect q^\top\vect k_i.
\end{equation}
然后对所有 $s_i$ 做 Softmax，得到读取各个 Value $\vect v_i$ 的权重。

\begin{center}
\begin{tikzpicture}[>=Latex,node distance=7mm,
 n/.style={draw=linegray,fill=paper,rounded corners=3pt,minimum width=107mm,minimum height=10mm,align=center,font=\small}]
\node[n] (q) {Query 和每个 Key 做点积 $\longrightarrow$ 匹配分数};
\node[n,below=of q] (s) {Softmax $\longrightarrow$ 一组非负、和为 1 的权重};
\node[n,below=of s] (v) {权重 $\times$ 各个 Value，再相加 $\longrightarrow$ 输出};
\draw[->,thick,blue] (q)--(s);
\draw[->,thick,blue] (s)--(v);
\end{tikzpicture}
\end{center}

这就是\textbf{注意力机制的基本读出过程}。你已经能理解它的核心数学，不需要先知道联想记忆的研究历史。

\subsection{先走完 Attention 主线，另一条数学路线留作选读}

我们的第一个问题已经解决了：\textbf{当前位置如何根据相关性直接读取上下文}。
因此下一章不再重新从外积讲起，而是接着问：
\textbf{同样的读取规则，怎样写成一个矩阵公式？}

如果你喜欢从“存入一条 Key--Value”的角度重新理解这套数学，
可以在学完 Transformer 后阅读第六部分的\textbf{联想记忆选读专题}。
它有助于研究记忆干扰、外积和非线性，却不是理解
Self-Attention 的先修条件。

\subsection{为什么 Attention 不等于无限长记忆？}

一次标准全量注意力需要让每个 Query 和很多 Key 比较；对 $T$ 个 Token，分数矩阵可达 $T\times T$。因此当 $T$ 增长时，注意力本身也可能带来显著计算和存储成本。它\textbf{缩短了信息交互路径}，并非免费解决所有长期依赖问题。

\begin{conclusion}
从 RNN 到 Attention，问题从“如何把过去信息逐步传下去”，变成“当前位置如何直接从上下文挑选需要的信息”。\textbf{下一章把这个单次查询扩展为矩阵计算}，不需要先走联想记忆的支线。
\end{conclusion}
